\section{总结与延伸}

\subsection{核心要点回顾}

本讲（Flow Matching 2）的核心贡献在于建立了扩散模型与流模型之间的\textbf{完整数学桥梁}，并展示了如何利用这一桥梁来设计更高效的生成模型。主要结论包括：

\begin{enumerate}
    \item \textbf{数学等价}：在线性高斯条件路径下，扩散模型是流模型的特殊情况。速度场 $u_t$、分数函数 $\nabla \log p_t$、噪声预测 $\boldsymbol{\epsilon}_\theta$ 三者可以精确互转

    \item \textbf{三种视角统一}：SDE（扩散）、PF-ODE（概率流）、Flow ODE（流匹配）在数学上完全一致

    \item \textbf{Rectified Flow}：线性插值定义条件轨迹，训练目标为速度预测，形式极其简洁

    \item \textbf{轨迹弯曲问题}：随机配对导致条件轨迹交叉，边际轨迹弯曲，需要多步 ODE 求解

    \item \textbf{Reflow 微调}：通过使用模型自身生成的配对数据重新训练，可有效将轨迹拉直，实现少步生成

    \item \textbf{Optimal Transport 联系}：OT 配对可减少轨迹交叉，Reflow 隐式逼近 OT 映射

    \item \textbf{实际影响}：SD3、Flux 等当前最先进模型已采用 Flow Matching / Rectified Flow 作为核心技术
\end{enumerate}

\subsection{拓展阅读}

\begin{itemize}
    \item Lipman, Y., et al. \textit{Flow Matching for Generative Modeling}. ICLR 2023. {\small Flow Matching 的奠基论文}
    \item Liu, X., Gong, C., \& Liu, Q. \textit{Flow Straight and Fast: Learning to Generate and Transfer Data with Rectified Flow}. ICLR 2023. {\small Rectified Flow 和 Reflow 的原始论文}
    \item Esser, P., et al. \textit{Scaling Rectified Flow Transformers for High-Resolution Image Synthesis}. ICML 2024. {\small SD3 的技术报告，将 Rectified Flow 扩展到大规模图像生成}
    \item Tong, A., et al. \textit{Improving and Generalizing Flow-Based Generative Models with Minibatch Optimal Transport}. TMLR 2024. {\small Mini-batch OT 在 Flow Matching 中的应用}
    \item Albergo, M. S. \& Vanden-Eijnden, E. \textit{Building Normalizing Flows with Stochastic Interpolants}. ICLR 2023. {\small Stochastic Interpolant 视角，与 Flow Matching 密切相关}
    \item Song, Y., et al. \textit{Consistency Models}. ICML 2023. {\small 加速采样的竞争方案，可与 Rectified Flow 对比}
\end{itemize}

%% --- Fin del contenido --- %%

\end{document}
